% ===================== 附录：练习详解 =====================
\clearpage
\phantomsection\label{sec:answers}
\hdr{附录}{练习详解（P1--P20）}

每块先复述题干与来源，再展开。正文只给题干与考点，步骤一律在这里；做完一题再翻一题，比连着读五题有用。

\hdrsub{A}{§7.1 的四道思考题（P1--P4）}

\answerstart{P1}{课件 p8 思考题}{同一种弦，把弦绷紧、换一根更粗的弦，音高分别怎么变？}
\textbf{条件摘要}：弦长 $l$ 固定、两端固定；波速 $a=\sqrt{T/\rho}$；音高由振动频率决定。

\textbf{考点}：$T$ 与 $\rho$ 各自的作用方向相反——一个在分子、一个在分母。

\textbf{关键一步}：先把"音高"翻译成"频率"，再用 $f=a/\lambda$ 把频率与波速连起来；弦长与边界给定了 $\lambda$，所以频率只随 $a$ 变。

\textbf{详解}：

弦上波的频率与波速、波长满足

\[
f=\frac{a}{\lambda}.
\]

两端固定的弦，可能的波长由弦长决定（$\lambda$ 与 $l$ 成正比，比例由边界条件与模态决定，本书不展开）。既然 $\lambda$ 只由几何决定，频率就与波速成正比：

\[
f\propto a=\sqrt{\frac{T}{\rho}}.
\]

\begin{itemize}
\item \textbf{绷紧}（$T$ 增大，$\rho$ 不变）：$a$ 增大 → $f$ 增大 → 音调升高。这就是调音时拧紧弦的道理。
\item \textbf{换粗弦}（$\rho$ 增大，$T$ 不变）：$a$ 减小 → $f$ 减小 → 音调降低。这就是吉他的低音弦又粗又沉的原因。
\end{itemize}

\textbf{答案}：绷紧 → 波速变大 → 音调升高；换粗弦 → 线密度变大 → 波速变小 → 音调降低。

\textbf{自检}：量纲上 $a=\sqrt{T/\rho}$ 是速度 $\checkmark$；方向上"紧则高、粗则低"与日常经验一致 $\checkmark$；若两种改动同时进行（比如换弦后又重新调音），结论要以最终的 $T/\rho$ 为准。

\qbreak

\answerstart{P2}{课件 p13 提问}{杆的纵振动加单位长度外力 $F(x,t)$ 后方程怎么写？外力是自重（沿 $-x$）时又怎么写？}
\textbf{条件摘要}：杆长 $l$、截面积 $S$、体密度 $\rho$、杨氏模量 $E$；$a^2=E/\rho$；外力以"单位长度"给出。

\textbf{考点}：非齐次项是"单位质量的外力"，不是力本身。

\textbf{关键一步}：微元 $(x,x+\Delta x)$ 受的外力是 $F\Delta x$，而质量是 $\rho S\Delta x$；牛顿第二定律右边除的是 $\rho S\Delta x$。

\textbf{详解}：

微元的受力多一项 $F\Delta x$，于是

\[
(\rho S\Delta x)\frac{\partial^2u}{\partial t^2}
=ES\left[\left.\frac{\partial u}{\partial x}\right|_{x+\Delta x}-\left.\frac{\partial u}{\partial x}\right|_{x}\right]+F\Delta x .
\]

两边除 $\rho S\Delta x$、取极限：

\[
\frac{\partial^2u}{\partial t^2}-a^2\frac{\partial^2u}{\partial x^2}=\frac{F(x,t)}{\rho S}.
\]

若外力是自重，方向沿 $-x$，单位长度重力为 $-\rho Sg$，代入得

\[
\frac{\partial^2u}{\partial t^2}-a^2\frac{\partial^2u}{\partial x^2}=-g .
\]

右边是常数，与 $x,t$ 无关。

\textbf{答案}：一般情形 $u_{tt}-a^2u_{xx}=F/(\rho S)$；自重情形 $u_{tt}-a^2u_{xx}=-g$。

\textbf{自检}：右边量纲是加速度（$\unit{m/s^2}$）$\checkmark$；杆自由悬挂时静止解满足 $-a^2u_{xx}=-g$，即 $u_{xx}=g/a^2=g\rho/E$，这就是自重下的均匀应变率 $\rho g/E$ $\checkmark$（与材料力学结果一致）。

\qbreak

\answerstart{P3}{课件 p3 提问}{独立推导一维热传导方程，说明 $a^2$ 与 $f$ 各由哪些物理量组成。}
\textbf{条件摘要}：均匀细杆，截面积 $S$、体密度 $\rho$、比热容 $c$、热导率 $\kappa$、单位体积生热率 $F$；温度 $u(x,t)$。

\textbf{考点}：三个守恒量各写一句——热量存量、净流入、源项。

\textbf{关键一步}：净流入要写"从 $x$ 面流入减从 $x+\Delta x$ 面流出"，符号靠傅里叶定律 $\mathbf q=-\kappa\nabla u$ 里的负号保证。

\textbf{详解}：

\textbf{(1) 热量存量}：微元体积 $S\Delta x$，单位体积热容 $\rho c$，故热量 $Q=\rho c\,u\,S\Delta x$，

\[
\frac{\dd Q}{\dd t}=\rho c\,S\Delta x\,\frac{\partial u}{\partial t}.
\]

\textbf{(2) 净流入}：傅里叶定律一维形式 $q_x=-\kappa\,\partial_xu$。净流入

\[
S q_x(x)-S q_x(x+\Delta x)
\approx-S\,\partial_xq_x\,\Delta x
=\kappa S\,\partial_x^2u\,\Delta x .
\]

\textbf{(3) 源项}：$F S\Delta x$。

\textbf{(4) 能量守恒}：

\[
\rho c\,S\Delta x\,u_t=\kappa S\,u_{xx}\Delta x+F S\Delta x
\ \Longrightarrow\
u_t-\frac{\kappa}{\rho c}u_{xx}=\frac{F}{\rho c}.
\]

\textbf{答案}：

\[
u_t-a^2u_{xx}=f,
\qquad
a^2=\frac{\kappa}{\rho c},\quad f=\frac{F}{\rho c}.
\]

$a^2$ 是热扩散率（\unit{m^2/s}），$f$ 是单位体积的"温度生成率"（$\unit{K/s}$）。

\textbf{自检}：$a^2=\kappa/(\rho c)$ 的量纲：$\kappa$ 是 $\unit{W\cdot m^{-1}\cdot K^{-1}}$，$\rho c$ 是 $\unit{J\cdot m^{-3}\cdot K^{-1}}$，相除得 $\unit{m^2/s}$ $\checkmark$；三项都化到 $\unit{K/s}$ $\checkmark$。三维推广只需把 $u_{xx}$ 换成 $\nabla^2u$。

\qbreak

\answerstart{P4}{按课件 p14 编排的数字练习}{膜的面密度 $\rho=0.020\ \unit{kg\cdot m^{-2}}$、张力 $T=2.0\ \unit{N\cdot m^{-1}}$：(1) 求波速；(2) 加均匀单位面积外力 $f=0.50\ \unit{m\cdot s^{-2}}$ 后写方程；(3) 说明 $f$ 的量纲。}
\textbf{条件摘要}：膜方程 $u_{tt}-a^2\nabla^2u=f$，$a=\sqrt{T/\rho}$，$\nabla^2=\partial_x^2+\partial_y^2$。

\textbf{考点}：$T$ 与 $\rho$ 在二维下的量纲；非齐次项被 $\rho$ 除过。

\textbf{关键一步}：先把 $T/\rho$ 的量纲算出来，再开方。

\textbf{详解}：

\textbf{(1)} 波速

\[
a=\sqrt{\frac{T}{\rho}}=\sqrt{\frac{2.0}{0.020}}=\sqrt{100}=10\ \unit{m/s}.
\]

量纲核对：$T$ 为 $\unit{N\cdot m^{-1}}=\unit{kg\cdot s^{-2}}$，$\rho$ 为 $\unit{kg\cdot m^{-2}}$，商为 $\unit{m^2\cdot s^{-2}}$，开方为 $\unit{m\cdot s^{-1}}$ $\checkmark$。

\textbf{(2)} 加均匀外力后

\[
\frac{\partial^2u}{\partial t^2}-a^2\left(\frac{\partial^2u}{\partial x^2}+\frac{\partial^2u}{\partial y^2}\right)=0.50 .
\]

\textbf{(3)} 方程各项量纲必须一致：$a^2\nabla^2u$ 的量纲是「$u$ 的量纲 ÷ $\unit{s^2}$」，所以右边 $f$ 也必须是「$u$ 的量纲 ÷ $\unit{s^2}$」。若 $u$ 以米计，$f$ 的量纲就是 $\unit{m\cdot s^{-2}}$（加速度）。原因是外力本来以「单位面积力」 $\unit{N\cdot m^{-2}}$ 给出，除以面密度 $\unit{kg\cdot m^{-2}}$ 后正好变成加速度。

\textbf{答案}：(1) $a=10\ \unit{m/s}$；(2) $u_{tt}-a^2\nabla^2u=0.50$；(3) $f$ 已被 $\rho$ 除过，量纲为加速度。

\textbf{自检}：$10\ \unit{m/s}$ 的量级对一张膜合理（比空气中声速慢 30 倍）；$f=0.50$ 与 $a^2u_{xx}$ 同量纲，可以相加 $\checkmark$。



\hdrsub{B}{§7.4 的四道定解条件题（P5--P8，课件 Ch7 p19）}

\answerstart{P5}{课件 Ch7 p19 第 1 题}{长 $\ell$ 的杆，左端固定；右端受沿外法向、大小为 $F(t)=F_0\sin\omega t$ 的力。写出边界条件。}
\explain{先看两端各“给”了什么。左端写“fixed”，给的是位移被按死，属于第一类，条件是函数值等于零；右端写“a normal force acts outward”，给的是力，属于第二类，条件是法向导数等于力除以刚度。右端的外法向朝 $+x$，力又朝外，所以 $f_n=+F(t)$，代数上不用再翻号——这是四题里最顺的一道。}

\qstem{\en{The left end of a rod (whose length is $l$) is fixed, and a normal force $F(t)=F_0\sin\omega t$ acts outward on the right end of the rod.}\\[3pt]
\textbf{中}（p19 第 1 题）：长为 $l$ 的杆，左端固定，右端受一外向法向力 $F(t)=F_0\sin\omega t$。求边界条件。}

\subhdr{解答过程}

\begin{enumerate}
\item 左端 $x=0$：固定端，位移恒为零，得第一类齐次条件 $u(0,t)=0$。
\item 右端 $x=l$：受外力。外法向 $n$ 指向 $+x$，力沿 $+x$ 向外，故 $f_n=+F_0\sin\omega t$。
\item 代入 Hooke 定律 $\dfrac{f_n}{S}=E\dfrac{\partial u}{\partial n}$，并用右端 $\partial u/\partial n=\partial u/\partial x$：
\[
\frac{F_0\sin\omega t}{S}=E\left.\frac{\partial u}{\partial x}\right|_{x=l}.
\]
\item 两边乘 $S$，写成作业常用的形式 $ES\,u_x(l,t)=F_0\sin\omega t$。
\end{enumerate}

\answerbox{$u(0,t)=0$（左端固定）；$\;ES\,u_x(l,t)=F_0\sin\omega t$，其中 $E$ 为杨氏模量、$S$ 为横截面积。（也可写成 $u_x|_{x=l}=F_0\sin\omega t/(ES)$。）}
{Boundary conditions: $u(0,t)=0$ (the left end is fixed) and $ES\,u_x(l,t)=F_0\sin\omega t$ (an outward normal force acts on the right end), where $E$ is Young's modulus and $S$ the cross-sectional area.}

\warnbox{\textbf{自检}：令 $F_0=0$，右端变成 $\left.u_x\right|_{x=l}=0$，正是 p14 的 “Free!”（自由端）\ok{}。反例：写成 $u(l,t)=0$ 就把右端当成了固定端，与“受拉力”矛盾。}

\answerstart{P6}{课件 Ch7 p19 第 2 题}{长 $\ell$ 的弦两端固定；$t=0$ 时在 $x=h$ 处受一外力 $F_0$。写出初始条件。}
\explain{这题问的是初始条件，不是边界条件——两端固定只是背景。“波动方程对时间二阶，所以初值要两个”：初始位移、初始速度。弦原本静止，位移初值为零；t=0 的一点受力，本质是一个\textbf{冲量}，它把速度在某一点上“打”成一个 $\delta$。判断用冲量—动量定理，量纲会替你检验：只有把 $F_0$ 当冲量（N$\cdot$s），$(F_0/\rho)\delta$ 才是速度。}

\qstem{\en{A string, whose length is $l$, is fixed at its end points. At $t=0$, an external force $F_0$ acts on the point $x=h$.}\\[3pt]
\textbf{中}（p19 第 2 题）：长为 $l$ 的弦，两端固定；$t=0$ 时在 $x=h$ 处受一外力 $F_0$。求初始条件。}

\subhdr{解答过程}

\begin{enumerate}
\item 先把问题分开：两端固定属于边界条件（$u(0,t)=u(l,t)=0$），本题不问；要写的是 $t=0$ 的两个初始条件。
\item 初始位移：$t<0$ 时弦静止在平衡位置，故 $u(x,0)=0$。
\item 初始速度：取 $x=h$ 附近长 $\Delta x$ 的一小段，冲量 $F_0$ 使它获得动量 $\rho\,\Delta x\,u_t$，于是 $u_t=\dfrac{F_0}{\rho\,\Delta x}$；令 $\Delta x\to0$，该速度集中在一点，写成 $\delta$ 函数：
\[
u_t(x,0)=\frac{F_0}{\rho}\,\delta(x-h).
\]
\item 量纲核对：$\dfrac{F_0}{\rho}$ 的单位是 $\dfrac{\mathrm{N\cdot s}}{\mathrm{kg/m}}=\mathrm{m^2/s}$，再乘 $\delta$ 的 $\mathrm{m^{-1}}$ 得 $\mathrm{m/s}$，正是速度 \ok{}。
\end{enumerate}

\answerbox{$u(x,0)=0$；$\;u_t(x,0)=\dfrac{F_0}{\rho}\,\delta(x-h)$，其中 $\rho$ 为弦的线密度，$F_0$ 按冲量理解。}
{Initial conditions: $u(x,0)=0$ and $u_t(x,0)=\dfrac{F_0}{\rho}\delta(x-h)$, where $\rho$ is the linear density of the string and $F_0$ is the impulse.}

\warnbox{\textbf{一个必须说清的约定}：若把 $F_0$ 当成普通力（单位 N）而不是冲量，它就不属于初始条件，而属于方程右端的源项 $\dfrac{F_0}{\rho}\delta(x-h)$，此时初条件只剩 $u(x,0)=0$、$u_t(x,0)=0$。两种理解只差“力持续多久”，量纲一眼可分。}

\answerstart{P7}{课件 Ch7 p19 第 3 题}{长 $\ell$ 的杆两端各受一个大小 $F_0$ 的法向力，方向都指向杆内。写出边界条件。}
\explain{两端给的都是力，所以两端都是第二类；真正要小心的是方向。图里两个箭头都指向杆内，是压力。外法向在两端相反，“向内”在两端对应的空间方向也相反，于是最省事的写法是先统一写成 $f_n$（沿外法向的分量）：压力与外法向反向，$f_n=-F_0$，两端一样；最后换到 $\partial u/\partial x$ 时，左端才翻号。}

\qstem{\en{A normal force $F_0$ acts on the ends of the rod (whose length is $l$).}\\[3pt]
\textbf{中}（p19 第 3 题）：长为 $l$ 的杆，两端各受一法向力 $F_0$（图中箭头指向杆内，即两端受压）。求边界条件。}

\rodfig{$F_0\;\longrightarrow$}{$\longleftarrow\;F_0$}

\subhdr{解答过程}

\begin{enumerate}
\item 两端都只给力、不给位移，故两端都是第二类边界条件。
\item 取外法向：左端 $n$ 沿 $-x$，右端 $n$ 沿 $+x$。压力方向与 $n$ 相反，所以两端都有 $f_n=-F_0$。
\item 代入 $\dfrac{f_n}{S}=E\dfrac{\partial u}{\partial n}$，得两端统一的形式
\[
\left.\frac{\partial u}{\partial n}\right|_{\text{两端}}=-\frac{F_0}{ES}.
\]
\item 换成 $x$ 导数：左端 $\partial u/\partial n=-\partial u/\partial x$，故 $-\left.u_x\right|_{x=0}=-\dfrac{F_0}{ES}$，即 $\left.u_x\right|_{x=0}=+\dfrac{F_0}{ES}$；右端 $\partial u/\partial n=+\partial u/\partial x$，直接得 $\left.u_x\right|_{x=l}=-\dfrac{F_0}{ES}$。
\end{enumerate}

\answerbox{$\left.u_x\right|_{x=0}=+\dfrac{F_0}{ES}$，$\;\left.u_x\right|_{x=l}=-\dfrac{F_0}{ES}$；等价地，外法向形式 $\left.\dfrac{\partial u}{\partial n}\right|_{\text{两端}}=-\dfrac{F_0}{ES}$。}
{$\left.u_x\right|_{x=0}=+F_0/(ES)$ and $\left.u_x\right|_{x=l}=-F_0/(ES)$; equivalently, in terms of the outward normal, $\partial u/\partial n=-F_0/(ES)$ at both ends.}

\warnbox{\textbf{与 p14 对照}：p14 是“外力向外”，得左端 $-\frac{f}{ES}$、右端 $+\frac{f}{ES}$；本题力向内，整组变号。\textbf{不要}只翻一端——若只把右端改成负号、左端照抄 p14，就得到两端不同向的怪结果。}

\answerstart{P8}{课件 Ch7 p19 第 4 题}{长 $\ell$ 的杆，两端各有热流 $q_0$ 流入杆内。写出边界条件。}
\explain{给的是热流密度，第二类；判据是 Fourier 定律，它自带负号。箭头朝内是“流入”，流入意味着热流逆着外法向，所以 $q_n=-q_0$，两端一样。换成 $\partial u/\partial x$ 后是“左负右正”，与题 3 正好相反——不是算错了，而是 Hooke 定律没有负号、Fourier 定律有负号。}

\qstem{\en{The heat flux into the ends of the rod (whose length is $l$) is $q_0$.}\\[3pt]
\textbf{中}（p19 第 4 题）：长为 $l$ 的杆，两端各流入热流密度 $q_0$。求边界条件。}

\rodfig{$q_0\;\longrightarrow$}{$\longleftarrow\;q_0$}

\subhdr{解答过程}

\begin{enumerate}
\item 未知量 $u$ 是温度，给定的是端面热流密度，故为第二类边界条件。
\item 取外法向 $n$，Fourier 定律 $q_n=-k\dfrac{\partial u}{\partial n}$；热流流入杆内即逆着 $n$，故 $q_n=-q_0$，两端相同。
\item 于是两端统一地有 $\dfrac{\partial u}{\partial n}=\dfrac{q_0}{k}$。
\item 换成 $x$ 导数：左端 $-\left.u_x\right|_{x=0}=\dfrac{q_0}{k}$，故 $\left.u_x\right|_{x=0}=-\dfrac{q_0}{k}$；右端 $\left.u_x\right|_{x=l}=\dfrac{q_0}{k}$ 直接得出。
\end{enumerate}

\answerbox{$\left.u_x\right|_{x=0}=-\dfrac{q_0}{k}$，$\;\left.u_x\right|_{x=l}=+\dfrac{q_0}{k}$，$k$ 为热导率；等价地 $-k\left.u_x\right|_{x=0}=q_0$、$-k\left.u_x\right|_{x=l}=-q_0$（后一式读作：$+x$ 方向净热流为 $-q_0$，真实热流沿 $-x$ 流入）。}
{$\left.u_x\right|_{x=0}=-q_0/k$ and $\left.u_x\right|_{x=l}=+q_0/k$ with $k$ the thermal conductivity; equivalently the inward heat flux at both ends is $q_0$.}

\warnbox{\textbf{自检}：两端吸热，端部比内部热，温度从左端往里降、从右端往里降，故左端 $u_x<0$、右端 $u_x>0$ \ok{}；$q_0=0$ 时两端绝热（$\left.u_x\right|=0$），回到 p15 的 adiabatic \ok{}。}

\hdrsub{C}{§8.1 的四道分离变量题（P9--P12，课件 Ch8 p13）}

%=====================================================================

\subhdr{四题对照：同一台机器换零件}

\tabtop
\begin{tabular}{@{}p{0.05\linewidth}p{0.17\linewidth}p{0.22\linewidth}p{0.24\linewidth}p{0.18\linewidth}@{}}
\thead{题} & \thead{方程类型} & \thead{特征值来自哪里} & \thead{定系数用的条件} & \thead{系数衰减} \\
\tabmidrule
1 & 波动 & $X(0)=X(\ell)=0$ & 初位移 $u(x,0)$ & $1/n^2$ \\
\tabmidrule
2 & 波动 & $X(0)=0,\ X'(\ell)=0$ & 初速度 $u_t(x,0)$ & $1/n^3$ \\
\tabmidrule
3 & 输运 & $X(0)=X(\ell)=0$ & 初温 $u(x,0)$ & $1/n^3$ \\
\tabmidrule
4 & 稳态 & $X(0)=X(a)=0$ & 非齐次边 $u(x,0)$ & $1/n$ \\
\end{tabular}
\tabbot

四道题把「分离变量」的四种组合各练了一遍：整数与半整数特征值、初位移与初速度、输运与稳态、一维与二维。考试里遇到新题，先在第 \pageref{sec:8111} 页的两张判定表里对号，再动笔。

\subhdr{交卷前的六条自检}

\begin{enumerate}[label=\textbf{(\arabic*)}]
  \item \textbf{边界。}把 $x=0$、$x=\ell$（或 $y=0$、$y=b$）代进最终级数，逐条看是否为零。这一步能抓住绝大部分符号错误。
  \item \textbf{初值。}令 $t=0$，看级数是否就是你给的初值（类型对不对：$\cos$ 还是 $\sin$，指数在不在）。
  \item \textbf{量纲。}系数乘上特征函数后必须与 $u$ 同量纲；系数里出现不该有的量纲，多半是漏除了一个 $\omega_n$ 或一个 $k_n$。
  \item \textbf{衰减。}光滑的初值对应快速衰减的系数；如果算出 $1/n$ 却预期 $1/n^3$，回去查积分或边界。
  \item \textbf{奇偶。}对称的初值不激发反对称模态；如果答案里本该消失的项没消失，先怀疑系数算错而不是怀疑对称性。
  \item \textbf{极限。}输运问题看 $t\to\infty$ 是否趋于稳态；Laplace 问题看 $A\to0$、$b\to\infty$ 等极限是否合理。
\end{enumerate}

%=====================================================================

\answerstart{P9}{课件 Ch8 p13 第 1 题}{长 $\ell$ 的弦两端固定，初始位移为峰在中点的三角波（斜率 $h$，峰值 $h\ell/2$），初始速度为零。求 $u(x,t)$。}
%=====================================================================

\qstem{（课件 p13 第 1 题）长 $\ell$ 的弦，两端固定。初始位移为
$\varphi(x)=hx\ (0<x<\ell/2)$、$\varphi(x)=h(\ell-x)\ (\ell/2<x<\ell)$，初始速度为零，$h$ 为常数。求 $u(x,t)$。}

\begin{figbox}
\begin{tikzpicture}[x=1cm,y=1cm,>=Stealth]
  \draw[->,inkgray] (-0.35,0)--(5.05,0) node[right,font=\small]{$x$};
  \draw[->,inkgray] (0,-0.30)--(0,2.35) node[above,font=\small]{$\varphi$};
  \draw[very thick,violet] (0,0)--(2.2,1.8)--(4.4,0);
  \draw[dashed,petal] (2.2,0)--(2.2,1.8);
  \node[below,font=\small] at (2.2,0) {$\ell/2$};
  \node[below,font=\small] at (4.4,0) {$\ell$};
  \node[left,font=\small] at (2.2,1.8) {$h\ell/2$};
  \node[font=\small,violetdeep] at (1.0,0.48) {$hx$};
  \node[font=\small,violetdeep] at (3.25,0.34) {$h(\ell-x)$};
\end{tikzpicture}
\par\smallskip{\small\color{inkgray}图 2-1\quad 初位移是底边为 $\ell$、峰在中点的三角波；$h$ 是左半段的斜率，峰值是 $h\ell/2$。}
\end{figbox}

\step{第 1 步：写定解问题}

弦的横振动满足 $u_{tt}=a^2u_{xx}$，其中 $a^2=T/\rho$（张力除以线密度）。两端固定是第一类齐次条件；初始位移给定、初始速度为零。于是

\[
\begin{cases}
u_{tt}=a^2u_{xx}, & 0<x<\ell,\ t>0,\\[2pt]
u(0,t)=u(\ell,t)=0, & t\ge 0,\\[2pt]
u(x,0)=\varphi(x),\quad u_t(x,0)=0, & 0\le x\le\ell .
\end{cases}
\]

\textbf{先记下量纲。}$\varphi=hx$ 说明 $[h]=[u]/L$，峰值是 $h\ell/2$，量纲为长度。最后核对系数时，$[C_n]$ 也必须是长度——这是检查「有没有把 $h$ 当峰值」的最快办法。

\step{第 2 步：分离变量}

设 $u(x,t)=X(x)T(t)$ 且不恒为零，代入方程得 $XT''=a^2X''T$，两边除以 $a^2XT$：

\[
\frac{T''}{a^2T}=\frac{X''}{X}=-\lambda .
\]

左边只含 $t$、右边只含 $x$，要恒等就必须都等于常数 $-\lambda$。于是得到一对常微分方程

\[
X''+\lambda X=0,
\qquad
T''+a^2\lambda T=0 .
\]

\textbf{为什么 $\lambda>0$。}这一步不能跳，因为「$\lambda$ 必须为正」是边界条件逼出来的，不是假设。

\begin{itemize}
  \item 若 $\lambda=-\kappa^2<0$：$X=A\cosh\kappa x+B\sinh\kappa x$。$X(0)=0$ 给出 $A=0$；再由 $X(\ell)=B\sinh\kappa\ell=0$ 得 $B=0$（因为 $\sinh\kappa\ell\ne0$）。只有零解，排除。
  \item 若 $\lambda=0$：$X=Ax+B$。$X(0)=0$ 给出 $B=0$，$X(\ell)=A\ell=0$ 给出 $A=0$。仍只有零解，排除。（$\lambda=0$ 只在两端都给「导数为零」时才可能，那时特征函数是常数。）
  \item 若 $\lambda=k^2>0$：$X=A\cos kx+B\sin kx$。$X(0)=0$ 给出 $A=0$；$X(\ell)=B\sin k\ell=0$。要非零解必须 $\sin k\ell=0$，即 $k\ell=n\pi$。
\end{itemize}

\step{第 3 步：解特征值问题}

由上一步，

\[
\lambda_n=\left(\frac{n\pi}{\ell}\right)^2,
\qquad
X_n(x)=\sin\frac{n\pi x}{\ell},
\qquad n=1,2,3,\dots
\]

$n$ 从 $1$ 开始：$n=0$ 给出 $X\equiv0$，不是解。这一步的结论可以和第 \pageref{sec:8111} 页的§8.1.2 的表 第一行直接对照。

\step{第 4 步：时间因子与初始速度}

把 $\lambda_n$ 代入时间方程：

\[
T_n''+\omega_n^2T_n=0,
\qquad
\omega_n=\frac{n\pi a}{\ell},
\]

其通解为 $T_n=C_n\cos\omega_nt+D_n\sin\omega_nt$。叠加起来

\[
u(x,t)=\sum_{n=1}^{\infty}\left(C_n\cos\omega_nt+D_n\sin\omega_nt\right)\sin\frac{n\pi x}{\ell}.
\]

现在用初始速度：对 $t$ 求导，

\[
u_t(x,t)=\sum_{n=1}^{\infty}\omega_n\left(-C_n\sin\omega_nt+D_n\cos\omega_nt\right)\sin\frac{n\pi x}{\ell}.
\]

令 $t=0$，得 $\sum_{n\ge1}\omega_nD_n\sin\frac{n\pi x}{\ell}=0$。两边乘 $\sin\frac{m\pi x}{\ell}$ 并在 $[0,\ell]$ 上积分，由正交性得 $\omega_mD_m\cdot\frac{\ell}{2}=0$，故每个 $D_m=0$。

于是级数里只剩余弦项：

\[
u(x,t)=\sum_{n=1}^{\infty}C_n\cos\frac{n\pi at}{\ell}\,\sin\frac{n\pi x}{\ell}.
\]

\textbf{物理读法。}余弦表示每个模态在 $t=0$ 时处在最大位移处、速度为零——这正是「把弦拉成三角波后松手」的初始状态。

\step{第 5 步：用初位移定系数}

令 $t=0$：

\[
\sum_{n=1}^{\infty}C_n\sin\frac{n\pi x}{\ell}=\varphi(x).
\]

两边乘 $\sin\frac{m\pi x}{\ell}$ 并从 $0$ 积到 $\ell$，右边只有 $n=m$ 项留下：

\[
C_n=\frac{2}{\ell}\int_0^{\ell}\varphi(x)\sin\frac{n\pi x}{\ell}\dd x
=\frac{2}{\ell}\left[\int_0^{\ell/2}hx\sin\beta x\dd x+\int_{\ell/2}^{\ell}h(\ell-x)\sin\beta x\dd x\right],
\qquad \beta=\frac{n\pi}{\ell}.
\]

两个积分都要分部积分。先备好两个原函数：

\[
\int x\sin\beta x\dd x=-\frac{x\cos\beta x}{\beta}+\frac{\sin\beta x}{\beta^2},
\qquad
\int(\ell-x)\sin\beta x\dd x=-\frac{(\ell-x)\cos\beta x}{\beta}-\frac{\sin\beta x}{\beta^2}.
\]

\textbf{左半段。}下限 $x=0$ 处两项都是零，只剩上限 $x=\ell/2$：

\[
I_1=\int_0^{\ell/2}hx\sin\beta x\dd x
=h\left[-\frac{(\ell/2)\cos(\beta\ell/2)}{\beta}+\frac{\sin(\beta\ell/2)}{\beta^2}\right].
\]

\textbf{右半段。}上限 $x=\ell$ 处 $\ell-x=0$ 且 $\sin\beta\ell=\sin n\pi=0$，整项为零；于是

\[
\begin{aligned}
I_2&=\int_{\ell/2}^{\ell}h(\ell-x)\sin\beta x\dd x\\
&=0-h\left[-\frac{(\ell/2)\cos(\beta\ell/2)}{\beta}-\frac{\sin(\beta\ell/2)}{\beta^2}\right]\\
&=h\left[\frac{(\ell/2)\cos(\beta\ell/2)}{\beta}+\frac{\sin(\beta\ell/2)}{\beta^2}\right].
\end{aligned}
\]

两段相加，含 $\cos(\beta\ell/2)$ 的项正好抵消：

\[
I_1+I_2=\frac{2h}{\beta^2}\sin\frac{\beta\ell}{2}
=\frac{2h\ell^2}{n^2\pi^2}\sin\frac{n\pi}{2}.
\]

代回 $C_n$ 的定义：

\[
C_n=\frac{2}{\ell}\cdot\frac{2h\ell^2}{n^2\pi^2}\sin\frac{n\pi}{2}
=\boxed{\ \frac{4h\ell}{n^2\pi^2}\sin\frac{n\pi}{2}\ }.
\]

\textbf{分奇偶写开。}当 $n=2k$ 时 $\sin k\pi=0$；当 $n=2k+1$ 时 $\sin\frac{(2k+1)\pi}{2}=(-1)^k$。所以

\[
C_{2k}=0,
\qquad
C_{2k+1}=\frac{4h\ell}{(2k+1)^2\pi^2}(-1)^k,
\qquad k=0,1,2,\dots
\]

\textbf{数值感受。}取 $\ell=1$、$h=1$（峰值 $0.5$）：$C_1=4/\pi^2\approx0.405$，$C_3=-4/(9\pi^2)\approx-0.045$，$C_5=4/(25\pi^2)\approx0.016$。三项相加在 $x=1/2$ 处给出 $0.405-0.045+0.016\approx0.376$，还差一点，正好说明收敛不快（$1/n^2$），需要更多项才能把尖角堆出来。

\step{第 6 步：最终答案}

\[
\boxed{\;
u(x,t)=\frac{4h\ell}{\pi^2}\sum_{k=0}^{\infty}\frac{(-1)^k}{(2k+1)^2}
\cos\frac{(2k+1)\pi at}{\ell}\,
\sin\frac{(2k+1)\pi x}{\ell}
\;}
\]

等价的紧凑写法是

\[
u(x,t)=\sum_{n=1}^{\infty}\frac{4h\ell}{n^2\pi^2}\sin\frac{n\pi}{2}\,
\cos\frac{n\pi at}{\ell}\,\sin\frac{n\pi x}{\ell},
\]

两种写法完全一样，考试里写哪一种都行，但要保证 $n$ 的取值说清楚（第二种里偶数项自动为零，不必额外声明）。

\step{第 7 步：自检}

\begin{itemize}
  \item \textbf{边界。}每个 $\sin\frac{n\pi x}{\ell}$ 在 $x=0$ 与 $x=\ell$ 处都为零，所以 $u(0,t)=u(\ell,t)=0$ 对任意 $t$ 成立。$\checkmark$
  \item \textbf{初速度。}级数对 $t$ 求导后每项都带 $\sin\omega_nt$，$t=0$ 时为零。$\checkmark$
  \item \textbf{初位移。}$C_n$ 正是 $\varphi$ 的傅里叶正弦系数，$t=0$ 时级数回到三角波。用 400 项截断在 $[0,\ell]$ 上比较，最大偏差 $5.1\times10^{-4}$（峰值 $0.5$），属于尖角处的截断余项，不是算错。
  \item \textbf{量纲。}$[C_n]=[h][\ell]$，与位移同量纲。$\checkmark$
  \item \textbf{为什么偶数项为零。}三角波关于 $x=\ell/2$ 对称，而 $n$ 为偶数时 $\sin\frac{n\pi x}{\ell}$ 关于中点反对称：对称的初值不能激发反对称的模态。也可以从节点看：偶数模态在 $x=\ell/2$ 处恰好为零，无法表现那里的尖角。
  \item \textbf{峰值换写。}若记峰值 $u_0=h\ell/2$，则 $C_n=\frac{8u_0}{n^2\pi^2}\sin\frac{n\pi}{2}$，与常见教材一致。用 $h$ 写时千万别漏 $\ell$，否则量纲就不对了。
  \item \textbf{收敛性。}位移级数按 $1/n^2$ 绝对一致收敛；逐项求导一次得到的 $u_t$ 按 $1/n$ 收敛，在尖角处只是条件收敛——这解释了为什么三角初位移的「速度场」在 $x=\ell/2$ 附近不好惹。
\end{itemize}

%=====================================================================

\answerstart{P10}{课件 Ch8 p13 第 2 题}{长 $\ell$ 的杆左端固定、右端自由；初始位移为零，初始速度 $\psi(x)=bx/\ell$。求 $u(x,t)$。}
%=====================================================================

\qstem{（课件 p13 第 2 题）长 $\ell$ 的杆，左端固定、右端自由。初始速度 $\psi(x)=bx/\ell$，初始位移为零，$b$ 为常数。求 $u(x,t)$。}

\step{第 1 步：写定解问题}

杆的纵振动同样满足 $u_{tt}=a^2u_{xx}$，其中 $a=\sqrt{E/\rho}$（杨氏模量除以密度）。左端固定是第一类齐次；右端自由表示端面不受力，应变 $\partial u/\partial x$ 为零，是第二类齐次：

\[
\begin{cases}
u_{tt}=a^2u_{xx}, & 0<x<\ell,\ t>0,\\[2pt]
u(0,t)=0,\quad u_x(\ell,t)=0, & t\ge0,\\[2pt]
u(x,0)=0,\quad u_t(x,0)=\dfrac{bx}{\ell}, & 0\le x\le\ell .
\end{cases}
\]

\textbf{与作业 1 的两处差别。}一是右端条件从「位移为零」换成「导数为零」，特征值会变；二是初值给的是速度而不是位移，定系数的路径不同。这两处差别恰好是本题的两个考点。

\textbf{量纲。}$\psi=bx/\ell$ 是速度，$[b]=L/T$；最后 $[A_n]=[b\ell/a]=L$，仍是长度。

\step{第 2 步：分离变量与特征值问题}

与作业 1 完全一样地得到 $X''+\lambda X=0$、$T''+a^2\lambda T=0$。$\lambda<0$ 与 $\lambda=0$ 的排除过程也相同（$\lambda=0$ 时 $X=Ax+B$，$X(0)=0$ 给 $B=0$，$X'(\ell)=A=0$，仍只有零解）。所以只讨论 $\lambda=k^2>0$：

\[
X=A\cos kx+B\sin kx .
\]

左端：$X(0)=A=0$，于是 $X=B\sin kx$。

右端：$X'(\ell)=Bk\cos k\ell=0$。要非零解，必须

\[
\cos k\ell=0
\quad\Longrightarrow\quad
k_n\ell=\left(n+\frac12\right)\pi
\quad\Longrightarrow\quad
k_n=\frac{(2n+1)\pi}{2\ell},
\qquad n=0,1,2,\dots
\]

这里 $n$ 从 $0$ 开始，因为 $k_0=\pi/(2\ell)$ 本身就是一个允许的模态（$\cos\frac{\pi}{2}=0$）。于是

\[
\lambda_n=\left(\frac{(2n+1)\pi}{2\ell}\right)^2,
\qquad
X_n(x)=\sin\frac{(2n+1)\pi x}{2\ell}.
\]

\textbf{一条容易记错的边界。}如果右端也是固定（$X(\ell)=0$），答案就退回 $\sin\frac{n\pi x}{\ell}$，与作业 1 相同；正因为右端是「自由」，才出现半整数。做题时先看两端是「值」还是「导数」，再决定用哪一组。

\step{第 3 步：时间因子与初始位移}

时间方程是 $T_n''+\omega_n^2T_n=0$，$\omega_n=ak_n=\frac{(2n+1)\pi a}{2\ell}$，通解为

\[
T_n=A_n\cos\omega_nt+B_n\sin\omega_nt .
\]

初始位移为零：$u(x,0)=0$ 要求每个 $T_n(0)=A_n=0$（同样用正交性逐项判断）。所以

\[
u(x,t)=\sum_{n=0}^{\infty}B_n\sin\omega_nt\,\sin\frac{(2n+1)\pi x}{2\ell}.
\]

为了和最终答案的记号一致，下面把这个待定系数记作 $A_n$（它乘在 $\sin\omega_nt$ 上，不是 $\cos$ 的系数）：

\[
u(x,t)=\sum_{n=0}^{\infty}A_n\sin\omega_nt\,\sin\frac{(2n+1)\pi x}{2\ell}.
\]

\step{第 4 步：用初速度定系数}

对 $t$ 求导并令 $t=0$：

\[
u_t(x,0)=\sum_{n=0}^{\infty}A_n\omega_n\sin\frac{(2n+1)\pi x}{2\ell}=\frac{bx}{\ell}.
\]

两边乘 $\sin\frac{(2m+1)\pi x}{2\ell}$ 并在 $[0,\ell]$ 上积分。半整数特征函数同样正交，模方仍为 $\ell/2$：

\[
A_n\omega_n=\frac{2}{\ell}\int_0^{\ell}\frac{bx}{\ell}\sin\frac{(2n+1)\pi x}{2\ell}\dd x
=\frac{2b}{\ell^2}\int_0^{\ell}x\sin k_nx\dd x .
\]

\textbf{算积分。}记 $k=k_n$，由分部积分

\[
\int_0^{\ell}x\sin kx\dd x
=\left[-\frac{x\cos kx}{k}+\frac{\sin kx}{k^2}\right]_0^{\ell}
=-\frac{\ell\cos k\ell}{k}+\frac{\sin k\ell}{k^2}.
\]

这里两个三角函数值都由 $k\ell=(n+\tfrac12)\pi$ 决定：

\[
\cos k\ell=0,
\qquad
\sin k\ell=\sin\left(n\pi+\frac{\pi}{2}\right)=(-1)^n .
\]

所以积分值只剩第二项：

\[
\int_0^{\ell}x\sin k_nx\dd x=\frac{(-1)^n}{k_n^2}
=\frac{4\ell^2(-1)^n}{(2n+1)^2\pi^2}.
\]

代回：

\[
A_n\omega_n=\frac{2b}{\ell^2}\cdot\frac{4\ell^2(-1)^n}{(2n+1)^2\pi^2}
=\frac{8b(-1)^n}{(2n+1)^2\pi^2}.
\]

\textbf{别忘了再除 $\omega_n$。}$\omega_n=\frac{(2n+1)\pi a}{2\ell}$，所以

\[
A_n=\frac{8b(-1)^n}{(2n+1)^2\pi^2}\cdot\frac{2\ell}{(2n+1)\pi a}
=\boxed{\ \frac{16b\ell\,(-1)^n}{(2n+1)^3\pi^3a}\ }.
\]

这一步是本题唯一容易翻车的地方：如果不除 $\omega_n$，直接写 $A_n=\frac{8b(-1)^n}{(2n+1)^2\pi^2}$，量纲立刻不对（速度量纲），而且 $t$ 的量纲也进不去。

\step{第 5 步：最终答案}

\[
\boxed{\;
u(x,t)=\sum_{n=0}^{\infty}\frac{16b\ell\,(-1)^n}{(2n+1)^3\pi^3a}
\sin\frac{(2n+1)\pi at}{2\ell}\,
\sin\frac{(2n+1)\pi x}{2\ell}
\;}
\]

展开前三项看看它长什么样：

\[
u=\frac{16b\ell}{\pi^3a}\left[
\sin\frac{\pi at}{2\ell}\sin\frac{\pi x}{2\ell}
-\frac{1}{27}\sin\frac{3\pi at}{2\ell}\sin\frac{3\pi x}{2\ell}
+\frac{1}{125}\sin\frac{5\pi at}{2\ell}\sin\frac{5\pi x}{2\ell}-\cdots
\right].
\]

\textbf{数值感受。}首项系数 $\frac{16}{\pi^3}\approx0.516$，第二项只有它的 $1/27\approx3.7\%$。所以 $t$ 稍大以后，杆的振动几乎就是「四分之一波长」那个基模在晃。

\step{第 6 步：自检}

\begin{itemize}
  \item \textbf{初始位移。}每项都带 $\sin\omega_nt$，$t=0$ 时全为零。$\checkmark$
  \item \textbf{初始速度。}求导后得到 $\sum A_n\omega_n\sin k_nx$，系数 $\frac{8b(-1)^n}{(2n+1)^2\pi^2}$ 正是 $\psi=bx/\ell$ 的半整数正弦系数。用 600 项截断比较，最大偏差 $3.4\times10^{-4}$。$\checkmark$
  \item \textbf{右端自由。}对 $x$ 求导，$\frac{\dd}{\dd x}\sin k_nx=k_n\cos k_nx$，在 $x=\ell$ 处 $\cos k_n\ell=0$，所以 $u_x(\ell,t)=0$。$\checkmark$
  \item \textbf{左端固定。}$\sin0=0$，$u(0,t)=0$。顺带核一下：$u_t(0,t)=0$ 与初速度在 $x=0$ 处为零（$\psi(0)=0$）是相容的。
  \item \textbf{量纲。}$[A_n]=[b][\ell]/([a])$，而 $[a]=L/T$、$[b]=L/T$，故 $[A_n]=L$。$\checkmark$
  \item \textbf{系数为什么是 $1/n^3$。}$A_n\omega_n$ 按 $1/n^2$ 衰减：$\psi=bx/\ell$ 在 $x=\ell$ 处值不为零，按「偶延拓」看过去在端点有折角，所以展开系数只能到 $1/n^2$；再除以 $\omega_n\sim n$，就得到 $1/n^3$。
\end{itemize}

%=====================================================================

\answerstart{P11}{课件 Ch8 p13 第 3 题}{长 $\ell$ 的杆两端温度为零，初始温度 $u(x,0)=bx(\ell-x)/\ell^2$。求温度分布。}
%=====================================================================

\qstem{（课件 p13 第 3 题）长 $\ell$ 的杆，两端温度为零，初始温度分布 $u(x,0)=bx(\ell-x)/\ell^2$，$b$ 为常数。求温度分布 $u(x,t)$。}

\step{第 1 步：写定解问题}

一维热传导 $u_t=a^2u_{xx}$，$a^2=k/(c\rho)$（热导率除以比热与密度之积）。两端零温是第一类齐次；初始温度给定：

\[
\begin{cases}
u_t=a^2u_{xx}, & 0<x<\ell,\ t>0,\\[2pt]
u(0,t)=u(\ell,t)=0, & t\ge0,\\[2pt]
u(x,0)=\dfrac{bx(\ell-x)}{\ell^2}, & 0\le x\le\ell .
\end{cases}
\]

\textbf{量纲。}$x(\ell-x)/\ell^2$ 是无量纲的，所以 $[b]=$ 温度。这一点对检查系数很有用：最终 $C_n$ 必须与 $b$ 同量纲。

\step{第 2 步：分离变量与特征值问题}

设 $u=XT$，代入得

\[
\frac{T'}{a^2T}=\frac{X''}{X}=-\lambda
\quad\Longrightarrow\quad
X''+\lambda X=0,\qquad T'+a^2\lambda T=0 .
\]

\textbf{注意与波动方程的区别。}时间方程是一阶的，只有 $T'=-a^2\lambda T$，所以时间因子只有指数、没有振荡。这是「输运」与「波动」最本质的差别：热传导没有惯性，不会来回晃。

边界条件与作业 1 完全相同（两端为零），所以特征值问题一字不改：

\[
\lambda_n=\left(\frac{n\pi}{\ell}\right)^2,
\qquad
X_n=\sin\frac{n\pi x}{\ell},
\qquad n=1,2,\dots
\]

\step{第 3 步：时间因子与叠加}

把 $\lambda_n$ 代入，$T_n'=-\frac{n^2\pi^2a^2}{\ell^2}T_n$，解为

\[
T_n(t)=\exp\left[-\left(\frac{n\pi a}{\ell}\right)^2t\right].
\]

叠加：

\[
u(x,t)=\sum_{n=1}^{\infty}C_n\exp\left[-\left(\frac{n\pi a}{\ell}\right)^2t\right]\sin\frac{n\pi x}{\ell}.
\]

每个模态都带一个负指数：$n$ 越大衰减越快，高频（细碎的温度起伏）先被抹平，剩下的是越来越光滑的低频形状。这就是热传导「抹平细节」的数学版本。

\step{第 4 步：用初温定系数}

令 $t=0$：

\[
u(x,0)=\sum_{n=1}^{\infty}C_n\sin\frac{n\pi x}{\ell}=\frac{bx(\ell-x)}{\ell^2}.
\]

由正交性，

\[
C_n=\frac{2}{\ell}\int_0^{\ell}\frac{bx(\ell-x)}{\ell^2}\sin\beta x\dd x
=\frac{2b}{\ell^3}\int_0^{\ell}x(\ell-x)\sin\beta x\dd x,
\qquad \beta=\frac{n\pi}{\ell}.
\]

\textbf{把积分拆开算。}记 $J=\int_0^{\ell}x(\ell-x)\sin\beta x\dd x=\ell\int_0^\ell x\sin\beta x\dd x-\int_0^\ell x^2\sin\beta x\dd x$。两个原函数是

\[
\int x\sin\beta x\dd x=-\frac{x\cos\beta x}{\beta}+\frac{\sin\beta x}{\beta^2},
\qquad
\int x^2\sin\beta x\dd x=-\frac{x^2\cos\beta x}{\beta}+\frac{2x\sin\beta x}{\beta^2}+\frac{2\cos\beta x}{\beta^3}.
\]

\textbf{上限 $x=\ell$。}利用 $\sin\beta\ell=\sin n\pi=0$、$\cos\beta\ell=(-1)^n$：

\[
\ell\left[-\frac{\ell\cos\beta\ell}{\beta}\right]-\left[-\frac{\ell^2\cos\beta\ell}{\beta}+\frac{2\cos\beta\ell}{\beta^3}\right]
=-\frac{2\cos\beta\ell}{\beta^3}.
\]

（含 $\ell^2\cos\beta\ell/\beta$ 的两项正好抵消。）

\textbf{下限 $x=0$。}只有 $x^2$ 那一个原函数的常数项不为零，取负号后贡献 $+\dfrac{2}{\beta^3}$。于是

\[
J=\frac{2}{\beta^3}\left[1-\cos\beta\ell\right]
=\frac{2\ell^3}{n^3\pi^3}\left[1-(-1)^n\right].
\]

代回：

\[
C_n=\frac{2b}{\ell^3}\cdot\frac{2\ell^3}{n^3\pi^3}\left[1-(-1)^n\right]
=\boxed{\ \frac{4b}{n^3\pi^3}\left[1-(-1)^n\right]\ }.
\]

分奇偶写开：

\[
C_{2k}=0,
\qquad
C_{2k+1}=\frac{8b}{(2k+1)^3\pi^3},
\qquad k=0,1,2,\dots
\]

\textbf{这里没有交错符号。}作业 1 的奇数项带 $(-1)^k$，本题不带——因为初温 $bx(\ell-x)/\ell^2$ 在中点两侧都为正（单峰），而三角波的奇数次模态在中点两侧符号相反，两者「对齐」的方式不同。这一点最容易在抄答案时抄错。

\step{第 5 步：最终答案}

\[
\boxed{\;
u(x,t)=\frac{8b}{\pi^3}\sum_{k=0}^{\infty}\frac{1}{(2k+1)^3}
\exp\left[-\left(\frac{(2k+1)\pi a}{\ell}\right)^2t\right]
\sin\frac{(2k+1)\pi x}{\ell}
\;}
\]

\step{第 6 步：自检}

\begin{itemize}
  \item \textbf{初温。}用 400 项截断在 $[0,\ell]$ 上比较 $u(x,0)$ 与 $bx(\ell-x)/\ell^2$，最大偏差 $1.5\times10^{-7}$——比作业 1 小三个数量级，因为初温光滑（两端值与导数都为零）。$\checkmark$
  \item \textbf{边界。}每个 $\sin\frac{n\pi x}{\ell}$ 在两端为零。$\checkmark$
  \item \textbf{长时间行为。}$t\to\infty$ 时每个指数都趋于零，$u\to0$：两端零温、热量不断流出，最终整根杆降到零度。$\checkmark$
  \item \textbf{为什么只剩奇数项。}初温关于 $x=\ell/2$ 对称，偶数模态关于中点反对称，对称的初值不激发反对称模态。与作业 1 是同一个理由。
  \item \textbf{平均温度的独立核对。}直接算初温的平均值：
  \[
  \frac{1}{\ell}\int_0^{\ell}\frac{bx(\ell-x)}{\ell^2}\dd x=\frac{b}{6}.
  \]
  再从级数算：$\frac{1}{\ell}\int_0^\ell\sin\frac{n\pi x}{\ell}\dd x=\frac{2}{n\pi}$（$n$ 奇），故平均值为
  \[
  \sum_{k\ge0}\frac{8b}{(2k+1)^3\pi^3}\cdot\frac{2}{(2k+1)\pi}
  =\frac{16b}{\pi^4}\sum_{k\ge0}\frac{1}{(2k+1)^4}
  =\frac{16b}{\pi^4}\cdot\frac{\pi^4}{96}=\frac{b}{6}.
  \]
  两条路给出同一个数，说明系数没有漏项、也没有多一个因子。这条自检推荐在考试里用，几十秒就能验完。
  \item \textbf{量纲与衰减。}$[C_n]=[b]$，正确；系数按 $1/n^3$ 衰减，级数一致收敛，且可逐项求导（热方程的解在 $t>0$ 后是解析的）。
\end{itemize}

%=====================================================================

\answerstart{P12}{课件 Ch8 p13 第 4 题}{矩形 $0<x<a$、$0<y<b$，$\nabla^2u=0$，边界为 $u(0,y)=u(a,y)=0$、$u(x,0)=A(1-x/a)$、$u(x,b)=0$。求 $u$。}
%=====================================================================

\qstem{（课件 p13 第 4 题）矩形域 $0<x<a$、$0<y<b$ 上解 Laplace 方程，边界条件为
$u(0,y)=u(a,y)=0$、$u(x,0)=A(1-x/a)$、$u(x,b)=0$。}

\begin{figbox}
\begin{tikzpicture}[x=1cm,y=1cm]
  \draw[very thick,violet] (0,0) rectangle (5.0,3.0);
  \node[font=\small] at (2.5,1.55) {$\nabla^2u=0$};
  \node[below,font=\small,violetdeep] at (2.5,-0.30) {$u(x,0)=A(1-x/a)$};
  \node[above,font=\small,violetdeep] at (2.5,3.32) {$u(x,b)=0$};
  \node[left,font=\small,violetdeep,rotate=90] at (-0.55,1.5) {$u(0,y)=0$};
  \node[right,font=\small,violetdeep,rotate=90] at (5.55,1.5) {$u(a,y)=0$};
  \node[font=\small] at (-0.45,-0.32) {$0$};
  \node[font=\small] at (5.05,-0.22) {$a$};
  \node[font=\small] at (-0.22,3.05) {$b$};
\end{tikzpicture}
\par\smallskip{\small\color{inkgray}图 5-1\quad 三条边齐次，只有下边 $y=0$ 非齐次。非齐次边在下面，这决定了 $y$ 方向用 $\sinh k(b-y)$ 而不是 $\sinh ky$。}
\end{figbox}

\step{第 1 步：写定解问题}

稳态、无内热源的温度场满足 Laplace 方程，四条边界给齐：

\[
\begin{cases}
u_{xx}+u_{yy}=0, & 0<x<a,\ 0<y<b,\\[2pt]
u(0,y)=u(a,y)=0, & 0\le y\le b,\\[2pt]
u(x,b)=0, & 0\le x\le a,\\[2pt]
u(x,0)=A\left(1-\dfrac{x}{a}\right), & 0\le x\le a .
\end{cases}
\]

注意边界条件的「三齐一非」结构：$x=0$、$x=a$、$y=b$ 三条边都是零，只有 $y=0$ 非齐次。这个结构决定了解的形式。

\step{第 2 步：决定沿哪个方向分离}

二维问题有两个方向，选哪个方向做特征值问题，取决于哪一对边界是齐次的：

\begin{itemize}
  \item $x$ 方向的两条边 $x=0$、$x=a$ 都给零——\textbf{齐次、成对}，适合做特征值问题；
  \item $y$ 方向只有 $y=b$ 给零，$y=0$ 是非齐次的——\textbf{不能}做特征值问题，只能解出通解后用两条边定常数。
\end{itemize}

所以令 $u=X(x)Y(y)$，$X$ 走特征值问题，$Y$ 走双曲函数，最后用 $y=0$ 的展开定系数。

\step{第 3 步：分离变量与 $X$ 的特征值}

代入 $u_{xx}+u_{yy}=0$：

\[
\frac{X''}{X}=-\frac{Y''}{Y}=-\lambda
\quad\Longrightarrow\quad
X''+\lambda X=0,\qquad Y''-\lambda Y=0 .
\]

$X$ 的边界条件是 $X(0)=X(a)=0$，与作业 1 完全同型（只是 $\ell$ 换成 $a$）：

\[
\lambda_n=\left(\frac{n\pi}{a}\right)^2,
\qquad
X_n(x)=\sin\frac{n\pi x}{a},
\qquad n=1,2,\dots
\]

\step{第 4 步：$Y$ 的方程与双曲函数}

把 $\lambda_n$ 代入：$Y''-k_n^2Y=0$，$k_n=\frac{n\pi}{a}$。通解是

\[
Y_n(y)=C_n\cosh k_ny+D_n\sinh k_ny .
\]

\textbf{用 $y=b$ 的条件。}$Y_n(b)=0$ 给出一个关系式，但更省事的做法是直接把它写成「在 $y=b$ 处为零」的形式：

\[
Y_n(y)=\sinh k_n(b-y).
\]

验证一下：$y=b$ 时 $\sinh0=0$；对 $y$ 求两次导，$\frac{\dd^2}{\dd y^2}\sinh k(b-y)=k^2\sinh k(b-y)$，所以满足 $Y''-k^2Y=0$。这两个条件它都自动满足，不需要再解方程组。

\textbf{为什么不用 $\sinh k_ny$。}$\sinh k_ny$ 满足 $y=0$ 处的条件（值为零），却满足不了 $y=b$ 处的条件；本题非齐次边在 $y=0$，齐次边在 $y=b$，所以要用 $\sinh k_n(b-y)$。如果像课件例 8.1.5 那样非齐次边在上边 $y=b$，就要换成 $\sinh k_ny$。\textbf{看哪条边是零，双曲函数就在哪条边上取零}——这是本题最实用的判断法。

\step{第 5 步：叠加与定系数}

把每个模态乘起来求和，并把分母 $\sinh k_nb$ 预先除掉（这样 $y=0$ 处的系数就是干净的 $A_n$）：

\[
u(x,y)=\sum_{n=1}^{\infty}A_n\,\frac{\sinh k_n(b-y)}{\sinh k_nb}\,\sin k_nx .
\]

先核一下这个写法没有破坏边界：$y=b$ 时分子为零；$x=0,a$ 时 $\sin k_nx=0$；只有 $y=0$ 处的条件还没用。

令 $y=0$，注意 $\sinh k_nb/\sinh k_nb=1$：

\[
\sum_{n=1}^{\infty}A_n\sin\frac{n\pi x}{a}=A\left(1-\frac{x}{a}\right).
\]

两边乘 $\sin\frac{m\pi x}{a}$ 并在 $[0,a]$ 上积分，模方为 $a/2$：

\[
A_n=\frac{2}{a}\int_0^aA\left(1-\frac{x}{a}\right)\sin\frac{n\pi x}{a}\dd x .
\]

\textbf{算积分。}令 $t=x/a$，$\dd x=a\dd t$：

\[
\int_0^a\left(1-\frac{x}{a}\right)\sin\frac{n\pi x}{a}\dd x
=a\int_0^1(1-t)\sin(n\pi t)\dd t .
\]

再分部积分（$u=1-t$，$\dd v=\sin n\pi t\dd t$）：

\[
\int_0^1(1-t)\sin n\pi t\dd t
=\left[-\frac{(1-t)\cos n\pi t}{n\pi}\right]_0^1-\int_0^1\frac{\cos n\pi t}{n\pi}\dd t
=\frac{1}{n\pi}-0=\frac{1}{n\pi}.
\]

（第二个积分 $\int_0^1\cos n\pi t\dd t=\frac{\sin n\pi}{n\pi}=0$。）

于是

\[
A_n=\frac{2A}{a}\cdot\frac{a}{n\pi}
=\boxed{\ \frac{2A}{n\pi}\ }.
\]

\step{第 6 步：最终答案}

\[
\boxed{\;
u(x,y)=\frac{2A}{\pi}\sum_{n=1}^{\infty}\frac{1}{n}\cdot
\frac{\sinh\frac{n\pi(b-y)}{a}}{\sinh\frac{n\pi b}{a}}\,
\sin\frac{n\pi x}{a}
\;}
\]

\textbf{读法。}每一项是「$x$ 方向的正弦驻波」乘「$y$ 方向从下往上衰减的双曲比例」。$n$ 越大，$\frac{\sinh k(b-y)}{\sinh kb}\approx e^{-ky}$ 衰减越快，所以高次项只在 $y=0$ 附近（下边界的热源附近）有影响；离下边越远，温度越接近零，符合「上边零温」的直觉。

\step{第 7 步：自检}

\begin{itemize}
  \item \textbf{上边。}$y=b$ 时 $\sinh0=0$，每一项为零，$u(x,b)=0$。数值上逐项代入，最大偏差为 $0$（机器精度）。$\checkmark$
  \item \textbf{两侧。}$x=0$ 与 $x=a$ 时 $\sin\frac{n\pi x}{a}=0$，$u=0$。$\checkmark$
  \item \textbf{下边。}$y=0$ 时级数就是 $A(1-x/a)$ 的正弦展开。用 2000 项在内部点比较，最大偏差 $2.0\times10^{-3}$；端点 $x=0$ 附近有吉布斯现象（把 $A(1-x/a)$ 按奇延拓看，$x=0$ 是真正的间断），所以端点不参与比较。$\checkmark$
  \item \textbf{方程。}逐项代入：$\partial_{xx}\sin k_nx\cdot\sinh=-\ k_n^2\sin k_nx\sinh$，$\partial_{yy}$ 给出 $+k_n^2\sin k_nx\sinh$，相加为零。$\checkmark$
  \item \textbf{最大值原理。}所有 $A_n>0$，$\sinh k_n(b-y)$ 随 $y$ 单调减，所以 $u$ 在 $y=0$ 最大、向 $y=b$ 单调降到零，$0\le u\le A$。与「内部无热源、边界温度非负」的定性结论一致。
  \item \textbf{量纲与极限。}$[A_n]=[A]$；$b\to\infty$ 时 $\sinh$ 比值趋于 $e^{-k_ny}$，解退化为半无限长条的情形；$A\to0$ 时解为零。
\end{itemize}

%=====================================================================

\subhdr{七个常见错误（按题号定位）}

%=====================================================================

\warnbox{\textbf{错误 1（作业 2 最常见）。}右端自由却仍用 $\lambda_n=(n\pi/\ell)^2$、$X_n=\sin\frac{n\pi x}{\ell}$。这两个式子只属于「两端都为零」。右端给的是 $u_x=0$，必须解 $\cos k\ell=0$，得到半整数。检验办法：把 $\sin\frac{n\pi x}{\ell}$ 的导数在 $x=\ell$ 处算一下，$\cos n\pi=(-1)^n\ne0$，立刻露馅。}

\warnbox{\textbf{错误 2（作业 1）。}把题目里的 $h$ 当成峰值，直接写 $C_n=\frac{8h}{n^2\pi^2}\sin\frac{n\pi}{2}$。课件里 $\varphi=hx$，峰值是 $h\ell/2$，正确系数是 $\frac{4h\ell}{n^2\pi^2}\sin\frac{n\pi}{2}$，差一个因子 $\ell/2$。量纲检查能立刻发现：$h$ 的量纲是「长度／长度」，不带 $\ell$ 的系数不是长度。}

\warnbox{\textbf{错误 3（作业 2）。}求出 $A_n\omega_n=\frac{8b(-1)^n}{(2n+1)^2\pi^2}$ 之后直接把它当 $A_n$，忘了除以 $\omega_n$。症状：最终答案里没有 $a$，量纲是速度而不是位移。}

\warnbox{\textbf{错误 4（作业 3）。}把 $C_n$ 写成带 $(-1)^n$ 的形式（照抄作业 1）。本题初温是单峰正函数，奇数项应当同号；$1-(-1)^n$ 展开只有 $+8b/(n^3\pi^3)$，没有交错因子。用平均温度 $\frac{b}{6}$ 核对最快。}

\warnbox{\textbf{错误 5（作业 3）。}指数写成 $\exp\left(-\frac{n\pi a}{\ell}t\right)$，漏了平方。正确的时间常数是 $\left(\frac{n\pi a}{\ell}\right)^2t$，量纲必须是「无量纲」：$a^2t/\ell^2$ 无量纲，$at/\ell$ 也无量纲，两者都能过这一关，所以要回到方程 $T'=-a^2\lambda_nT$ 去核，$\lambda_n$ 带平方。}

\warnbox{\textbf{错误 6（作业 4）。}$Y$ 用 $\sinh k_ny$，结果 $y=b$ 处不为零。正确的是 $\sinh k_n(b-y)$（非齐次边在下）或 $\sinh k_ny$（非齐次边在上）。口诀：\textbf{双曲函数在零边界那条边上取零}。}

\warnbox{\textbf{错误 7（作业 4）。}忘记除 $\sinh k_nb$，却仍用 $A_n=\frac{2A}{n\pi}$。这两个选择必须配套：不除分母时，$y=0$ 处的方程变成 $\sum A_n\sinh k_nb\sin k_nx=A(1-x/a)$，系数要相应改成 $\frac{2A}{n\pi\sinh k_nb}$。两种写法都对，混用就错。}

%=====================================================================

\hdrsub{D}{§8.2 的四道极坐标题（P13--P16，自编）}

%=====================================================================

\answerstart{P13}{自编，圆内 Dirichlet 问题}{单位圆盘 $r<1$，$\nabla^2u=0$，$u(1,\varphi)=\sin^2\varphi$。求 $u$ 及其直角坐标形式。}

\textbf{第 1 步（把边界值写成级数）。}用半角公式

\[
\sin^2\varphi=\frac{1-\cos2\varphi}{2}=\frac12-\frac12\cos2\varphi .
\]

这是一个只含常数与 $\cos2\varphi$ 的式子，正好是 $\{1,\cos m\varphi\}$ 的有限组合。

\textbf{第 2 步（圆内通解与筛项）。}圆内解必须是 \eqref{eq:disksol} 的形状（$\ln r$ 与 $r^{-m}$ 在原点发散，丢掉）：

\[
u(r,\varphi)=C_0+\sum_{m\ge1}r^m\big(A_m\cos m\varphi+B_m\sin m\varphi\big).
\]

在 $r=1$ 处令其等于边界值：

\[
C_0+\sum_{m\ge1}\big(A_m\cos m\varphi+B_m\sin m\varphi\big)=\frac12-\frac12\cos2\varphi .
\]

\textbf{第 3 步（逐阶比较）。}左边是 $\varphi$ 的 Fourier 级数，右边只含两项。由正交性（乘 $1$、$\cos2\varphi$ 后积分），

\[
C_0=\frac12,\qquad A_2=-\frac12,\qquad \text{其余 }A_m=B_m=0 .
\]

\textbf{第 4 步（写出解）。}

\[
u(r,\varphi)=\frac12-\frac12r^2\cos2\varphi .
\]

\textbf{第 5 步（化成直角坐标）。}由 $r^2\cos2\varphi=r^2(\cos^2\varphi-\sin^2\varphi)=x^2-y^2$，

\[
u(x,y)=\frac12-\frac{x^2-y^2}{2}=\frac12-\frac{x^2}{2}+\frac{y^2}{2}.
\]

\textbf{第 6 步（核对）。}直接求导：$u_{xx}=-1$，$u_{yy}=1$，故 $u_{xx}+u_{yy}=0$，满足方程。在 $r=1$ 上回到极坐标：$\frac12-\frac12\cos2\varphi=\sin^2\varphi$，满足边界条件。两条都对，解确定（唯一性见 \S8.2.6.4）。

\aside{这个解是“常数 + 一个二次调和多项式”。在极坐标通解里，二次项只有 $r^2\cos2\varphi=x^2-y^2$ 与 $r^2\sin2\varphi=2xy$ 两个独立形式；一次项是 $x=r\cos\varphi$ 与 $y=r\sin\varphi$，零次项是 $1$。本题用的是 $1$ 与 $x^2-y^2$。}

\answerstart{P14}{自编，环域 $m=0$ 分支}{$a<r<b$，$u(a)=0$，$u(b)=U_0$，求 $u$、$E_r$、$\sigma_a$ 与单位长度电容。}

\textbf{第 1 步（筛项）。}两条边界值都与 $\varphi$ 无关，是常数。由正交性，所有 $m\ge1$ 的系数为零（否则右边会出现 $\cos m\varphi$、$\sin m\varphi$ 项）。于是只剩 $m=0$ 的两个常数：

\[
u(r)=C_0+D_0\ln r .
\]

\textbf{第 2 步（定两个常数）。}

\[
\begin{cases}
C_0+D_0\ln a=0,\\[2pt]
C_0+D_0\ln b=U_0 .
\end{cases}
\]

两式相减：$D_0\ln\dfrac{b}{a}=U_0$，故 $D_0=\dfrac{U_0}{\ln(b/a)}$；代回第一式 $C_0=-D_0\ln a$。于是

\[
u(r)=U_0\,\frac{\ln(r/a)}{\ln(b/a)} .
\]

核对：$r=a$ 时 $\ln1=0$，$u=0$；$r=b$ 时 $\ln(b/a)/\ln(b/a)=1$，$u=U_0$。都对。

\textbf{第 3 步（场）。}

\[
E_r=-\frac{\dd u}{\dd r}=-\frac{U_0}{r\ln(b/a)},\qquad E_\varphi=0 .
\]

$E_r<0$（$U_0>0$ 时）表示场指向内筒：外筒电位高，场从外向内。这与“内筒带负电、外筒带正电”一致。

\textbf{第 4 步（内筒表面电荷）。}

\[
\sigma_a=\varepsilon_0E_r(a)=-\frac{\varepsilon_0U_0}{a\ln(b/a)},\qquad
Q'=2\pi a\,\sigma_a=-\frac{2\pi\varepsilon_0U_0}{\ln(b/a)} .
\]

\textbf{第 5 步（电容）。}

\[
C=\frac{|Q'|}{U_0}=\frac{2\pi\varepsilon_0}{\ln(b/a)} .
\]

\textbf{第 6 步（极限对照）。}当 $b=a+d$、$d\ll a$ 时，$\ln(b/a)=\ln(1+d/a)\approx d/a$，于是

\[
C\approx\frac{2\pi\varepsilon_0 a}{d}=\frac{\varepsilon_0\,(2\pi a)}{d},
\]

正是“面积为 $2\pi a$（单位长度）、间距为 $d$”的平行板电容。这说明 $\ln r$ 这一支不是数学游戏：它在大间距时给出对数衰减，在小间距时退化为平行板。

\answerstart{P15}{自编，半圆盘的固定端点条件}{$r<a,\ 0<\varphi<\pi$，$\nabla^2u=0$；$u(r,0)=u(r,\pi)=0$，$u(a,\varphi)=U_0$。求级数解。}

\textbf{第 1 步（角向问题换成固定端点）。}分离后仍是 $\Phi''+\lambda\Phi=0$，但 $\varphi$ 的区间是 $(0,\pi)$，条件不再是“周期”，而是两端为零：

\[
\Phi(0)=0,\qquad \Phi(\pi)=0 .
\]

\textbf{第 2 步（定谱）。}能量恒等式 \eqref{eq:rayleigh} 在这里照样成立（边界项 $\Phi\Phi'$ 在两端为零），所以 $\lambda\ge0$。$\lambda=0$ 时 $\Phi=C_1+C_2\varphi$，两端为零逼 $C_1=C_2=0$，舍去。$\lambda=k^2>0$ 时 $\Phi=A\cos k\varphi+B\sin k\varphi$；$\Phi(0)=0$ 给 $A=0$，$\Phi(\pi)=0$ 给 $B\sin k\pi=0$，故 $k=n$（正整数）。于是

\[
\lambda=n^2,\qquad \Phi_n(\varphi)=\sin n\varphi,\qquad n=1,2,3,\dots
\]

注意 $\cos n\varphi$ 被排除：它在 $\varphi=0$ 处等于 $1$，不可能为零。这就是“固定端点”与“周期”两种条件的差别——谱都是整数平方，但特征函数只取一半。

\textbf{第 3 步（径向解与筛项）。}$R=r^n$ 与 $r^{-n}$ 都满足欧拉方程；圆内 $r\to0$ 有界，丢掉 $r^{-n}$。于是

\[
u(r,\varphi)=\sum_{n\ge1}A_n\,r^n\sin n\varphi .
\]

\textbf{第 4 步（用 $r=a$ 的条件定系数）。}

\[
\sum_{n\ge1}A_n\,a^n\sin n\varphi=U_0,\qquad 0<\varphi<\pi .
\]

在 $[0,\pi]$ 上用正弦系的正交性 $\displaystyle\int_0^\pi\sin n\varphi\sin m\varphi\dd\varphi=\frac\pi2\delta_{mn}$，两边乘 $\sin n\varphi$ 积分：

\[
A_n\,a^n\cdot\frac\pi2=\int_0^\pi U_0\sin n\varphi\dd\varphi=\frac{U_0\,(1-\cos n\pi)}{n}.
\]

所以

\[
A_n\,a^n=\frac{2U_0}{n\pi}\big(1-(-1)^n\big)=
\begin{cases}
\dfrac{4U_0}{n\pi}, & n\text{ 为奇数},\\[6pt]
0, & n\text{ 为偶数}.
\end{cases}
\]

\textbf{第 5 步（写出级数）。}令 $n=2k+1$：

\begin{equation}
u(r,\varphi)=\frac{4U_0}{\pi}\sum_{k=0}^{\infty}\frac{1}{2k+1}\Big(\frac{r}{a}\Big)^{2k+1}\sin\big((2k+1)\varphi\big).
\label{eq:halfdisk}
\end{equation}

前三项：

\[
u=\frac{4U_0}{\pi}\left[\Big(\frac ra\Big)\sin\varphi+\frac13\Big(\frac ra\Big)^{3}\sin3\varphi+\frac15\Big(\frac ra\Big)^{5}\sin5\varphi+\cdots\right].
\]

\textbf{第 6 步（在边界上收敛到什么）。}在 $r=a$ 上，\eqref{eq:halfdisk} 是常数 $U_0$ 在 $(0,\pi)$ 上的正弦级数。由 Fourier 收敛定理，它在每点 $0<\varphi<\pi$ 收敛到 $U_0$；在两端 $\varphi=0,\pi$ 各项为零，级数收敛到 $0$，而那里题给的条件也是 $0$（$u(r,0)=u(r,\pi)=0$），所以端点上的值也与条件一致。

\textbf{第 7 步（两个角点）。}在角点 $(a,0)$ 与 $(a,\pi)$，圆弧上的边界值 $U_0$ 与直径上的边界值 $0$ 直接相接，题给的数据本身在这两点跳变。因此解在角点附近不可能连续地取到两边的值：靠近角点时级数会出现 Gibbs 现象（在 $r=a$ 附近振荡），而在区域内部 $r<a$ 一切正常。严格地说，这题的解满足“区域内部方程 + 去掉两个角点后的边界条件”，在角点上不存在连续的边界延拓。这是半圆域问题的固有性质，不是解法的问题。

\aside{若把圆盘换成四分之一圆，谱仍是 $n^2$ 但特征函数变成 $\sin(2n\varphi)$（因为在 $\varphi=\pi/2$ 处也要为零，$n$ 必须为偶数）。谱的“整数平方”来自区间长度 $\pi$ 与两端为零，换成别的角宽就会变成 $\lambda=(n\pi/\alpha)^2$——这与 \S8.2.3.4 里“直角坐标的谱依赖区域尺寸”是同一件事。}

\answerstart{P16}{自编，判断四个函数能否作圆内的解}{单位圆盘 $r<1$ 内，(a) $r^{-1}\cos\varphi$；(b) $\ln r$；(c) $r^2\sin2\varphi$；(d) $r^{1/2}\cos(\varphi/2)$，哪几个可以是 $\nabla^2u=0$ 的解？}

判据只有两条，都来自 \S8.2.6：\textbf{(i) 在圆内（不含 $r=0$ 的奇点）满足 $\nabla^2u=0$；(ii) 在 $r=0$ 处有界，且对 $\varphi$ 以 $2\pi$ 为周期（单值）。}题干已说四个都满足 (i)，所以只需查 (ii)。

\textbf{(a) $r^{-1}\cos\varphi$。}它是 $m=1$ 的 $r^{-m}$ 项，属于通解 \eqref{eq:general}，在 $r>0$ 上调和。但 $r\to0$ 时 $\dfrac{\cos\varphi}{r}$ 沿 $\varphi=0$ 方向趋于 $+\infty$、沿 $\varphi=\pi$ 方向趋于 $-\infty$，无界。\textbf{不能}作为圆内的解（它是圆外或环域里的偶极线场）。

\textbf{(b) $\ln r$。}这是 $m=0$ 的第二个解，$\nabla^2\ln r=0$（$r>0$）。但 $r\to0^+$ 时 $\ln r\to-\infty$，无界。\textbf{不能}作为圆内的解。（它的物理对应是原点处的一根线电荷；圆内若真有线电荷，方程右边就不是 $0$ 而是 $\delta$ 函数，问题性质就变了。）

\textbf{(c) $r^2\sin2\varphi$。}这是 $m=2$ 的 $r^m$ 项，属于 \eqref{eq:disksol}。$r\le1$ 时 $|r^2\sin2\varphi|\le1$，有界；$\sin2\varphi$ 以 $2\pi$ 为周期，单值。\textbf{可以}。写成直角坐标：$r^2\sin2\varphi=2r^2\sin\varphi\cos\varphi=2xy$，而 $\partial_x^2(2xy)+\partial_y^2(2xy)=0$，与判据一致。

\textbf{(d) $r^{1/2}\cos(\varphi/2)$。}它在 $r>0$ 上调和：记 $u=r^{1/2}\cos\frac\varphi2$，则

\[
u_{rr}=-\frac14r^{-3/2}\cos\frac\varphi2,\qquad
\frac{u_r}{r}=\frac12r^{-3/2}\cos\frac\varphi2,\qquad
\frac{u_{\varphi\varphi}}{r^2}=-\frac14r^{-3/2}\cos\frac\varphi2,
\]

三项相加为零。$r\to0$ 时 $r^{1/2}\to0$，也有界。但把 $\varphi$ 换成 $\varphi+2\pi$：

\[
r^{1/2}\cos\Big(\frac{\varphi+2\pi}{2}\Big)=r^{1/2}\cos\Big(\frac\varphi2+\pi\Big)=-r^{1/2}\cos\frac\varphi2 ,
\]

变号了。绕原点一圈回来，函数值变成相反数，不是单值函数；沿 $\varphi=\pi$（负 $x$ 轴）它从 $0$ 跳到 $0$ 但两侧符号相反，实际上是两支拼在一起。\textbf{不能}作为圆域上的解。它对应的物理情形是“割线问题”（例如一张沿负实轴剪开的平面），不属于本节的圆域问题。

\textbf{结论}：只有 (c) 可以。

\aside{这道题的用处是把“通解”当筛子用：(a)(b) 是通解里被有界性筛掉的项，(c) 是留下来的项，(d) 是连 $\varphi$ 的周期性都不满足的项。考试里常见的问法是“下列函数中哪个是圆内调和函数”，判据就是这两条，不要只查 $\nabla^2u=0$。}

\hdrsub{E}{§8.3 的四道柱坐标题（P17--P20，含课件 p6 作业）}

%=====================================================================

\answerstart{P17}{课件 p6 作业 5}{圆内 $r<\rho_0$，$\nabla^2u=0$；边界值分别为 $A\cos\varphi$ 与 $A+B\sin\varphi$。}

\textbf{第 1 步（圆内通解）。}由 \S8.2 的结论，圆内有界解是

\[
u(r,\varphi)=C_0+\sum_{m\ge1}r^m\big(a_m\cos m\varphi+b_m\sin m\varphi\big)
\]

（用 $a_m,b_m$ 以免与题给的常数 $A,B$ 混淆）。

\textbf{第 2 步（情形 (1)）。}边界值只含 $\cos\varphi$ 一项，故只有 $m=1$ 的余弦系数非零：

\[
a_1\rho_0=A,\qquad C_0=0,\qquad \text{其余为零},
\]

所以

\[
u(r,\varphi)=A\,\frac{r}{\rho_0}\cos\varphi=A\,\frac{x}{\rho_0}.
\]

\textbf{核对。}$\Delta u=0$（$x$ 是调和函数）；$r=\rho_0$ 时 $u=A\cos\varphi$，对。

\textbf{第 3 步（情形 (2)）。}边界值含常数与 $\sin\varphi$ 两项：

\[
C_0=A,\qquad b_1\rho_0=B,\qquad \text{其余为零},
\]

所以

\[
u(r,\varphi)=A+B\,\frac{r}{\rho_0}\sin\varphi=A+B\,\frac{y}{\rho_0}.
\]

\textbf{核对。}常数与 $y$ 都调和；$r=\rho_0$ 时回到 $A+B\sin\varphi$，对。两问都说明一件事：\textbf{边界值只含有限个谐波时，解只含那几项}，不需要写无穷级数。

\answerstart{P18}{课件 p6 作业 7}{环域 $\rho_1<r<\rho_2$，$\nabla^2u=0$；$u(\rho_1)=u_1\cos2\varphi$，$u(\rho_2)=u_2\sin\varphi$。}

\textbf{第 1 步（通解的形式）。}环域保留正负幂次：

\[
u(r,\varphi)=C_0+D_0\ln r+\sum_{m\ge1}\big(A_mr^m+C_mr^{-m}\big)\cos m\varphi+\sum_{m\ge1}\big(B_mr^m+D_mr^{-m}\big)\sin m\varphi .
\]

\textbf{第 2 步（筛项）。}两条边界值里只有 $\cos2\varphi$ 与 $\sin\varphi$，所以只保留 $m=2$ 的余弦与 $m=1$ 的正弦：

\[
u(r,\varphi)=\big(A_2r^2+C_2r^{-2}\big)\cos2\varphi+\big(B_1r+B_1'r^{-1}\big)\sin\varphi .
\]

（$C_0=D_0=0$：边界值没有常数项；其它谐波系数为零：正交性。）

\textbf{第 3 步（定 $\cos2\varphi$ 的系数）。}由 $r=\rho_1$ 与 $r=\rho_2$：

\[
\begin{cases}
A_2\rho_1^2+C_2\rho_1^{-2}=u_1,\\[2pt]
A_2\rho_2^2+C_2\rho_2^{-2}=0 .
\end{cases}
\]

第二式给 $C_2=-A_2\rho_2^4$；代入第一式：$A_2\big(\rho_1^2-\rho_2^4\rho_1^{-2}\big)=u_1$，即

\[
A_2=\frac{u_1\rho_1^2}{\rho_1^4-\rho_2^4},\qquad C_2=\frac{u_1\rho_1^2\rho_2^4}{\rho_2^4-\rho_1^4}.
\]

\textbf{第 4 步（定 $\sin\varphi$ 的系数）。}

\[
\begin{cases}
B_1\rho_1+B_1'\rho_1^{-1}=0,\\[2pt]
B_1\rho_2+B_1'\rho_2^{-1}=u_2,
\end{cases}
\quad\Longrightarrow\quad
B_1=\frac{u_2\rho_2}{\rho_2^2-\rho_1^2},\qquad B_1'=-\frac{u_2\rho_1^2\rho_2}{\rho_2^2-\rho_1^2}.
\]

\textbf{第 5 步（写出答案并核对）。}

\[
u=A_2\Big(r^2-\frac{\rho_2^4}{r^2}\Big)\cos2\varphi
+B_1\Big(r-\frac{\rho_1^2}{r}\Big)\sin\varphi ,
\]

其中 $A_2,B_1$ 如上。核对：$r=\rho_1$ 时第二项为零、第一项为 $A_2(\rho_1^2-\rho_2^4/\rho_1^2)\cos2\varphi=u_1\cos2\varphi$；$r=\rho_2$ 时第一项为零、第二项为 $B_1(\rho_2-\rho_1^2/\rho_2)\sin\varphi=u_2\sin\varphi$。两处都对。

\aside{本题最容易错的地方是把两个不同阶的谐波混在一个方程组里。它们的 $\varphi$ 依赖不同（$\cos2\varphi$ 与 $\sin\varphi$），由正交性各自独立，必须分开解两个 $2\times2$ 方程组。}

\answerstart{P19}{自编，Fourier--Bessel 展开}{$0<r<a$，$f(r)=1-r^2/a^2$，用 $J_0(x_n^{(0)}r/a)$ 展开。}

\textbf{第 1 步（系数公式）。}由 \eqref{eq:fb}，取 $m=0$、$k_{0n}=x_n^{(0)}/a$：

\[
A_n=\frac{2}{a^2J_1^2(x_n^{(0)})}\int_0^a\Big(1-\frac{r^2}{a^2}\Big)J_0\Big(\frac{x_n^{(0)}r}{a}\Big)r\dd r .
\]

\textbf{第 2 步（算积分）。}令 $t=x_n^{(0)}r/a$，则 $r=a t/x_n$，$r\dd r=(a/x_n)^2t\dd t$：

\[
\int_0^a\Big(1-\frac{r^2}{a^2}\Big)J_0\Big(\frac{x_nr}{a}\Big)r\dd r
=\Big(\frac{a}{x_n}\Big)^2\int_0^{x_n}\Big(1-\frac{t^2}{x_n^2}\Big)J_0(t)\,t\dd t .
\]

用 $\int_0^ztJ_0(t)\dd t=zJ_1(z)$、$\int_0^zt^2J_1(t)\dd t=z^2J_2(z)$ 与分部积分（$t^3J_0=\frac{\dd}{\dd t}[t^3J_1]-2t^2J_1$），得

\[
\int_0^{x_n}\Big(1-\frac{t^2}{x_n^2}\Big)J_0(t)t\dd t
=2J_2(x_n)
=\frac{4J_1(x_n)}{x_n}\quad(\text{用到 }J_0(x_n)=0,\ J_2=\tfrac{2}{x}J_1-J_0).
\]

\textbf{第 3 步（系数）。}把上式代回：

\[
A_n=\frac{2}{a^2J_1^2(x_n)}\cdot\Big(\frac{a}{x_n}\Big)^2\cdot\frac{4J_1(x_n)}{x_n}
=\frac{8}{x_n^3J_1(x_n)} .
\]

\textbf{第 4 步（前三个数值）。}用 $x_1=2.4048$、$J_1(x_1)=0.5192$；$x_2=5.5201$、$J_1(x_2)=-0.3403$；$x_3=8.6537$、$J_1(x_3)=0.2715$：

\begin{gather*}
A_1=\frac{8}{2.4048^3\times0.5192}=1.108,\qquad
A_2=\frac{8}{5.5201^3\times(-0.3403)}=-0.140,\\
A_3=\frac{8}{8.6537^3\times0.2715}=0.0455 .
\end{gather*}

系数按 $1/x_n^3$ 衰减，比 $f=1$ 那个例子快得多——因为 $1-r^2/a^2$ 在 $r=a$ 处等于零，没有端点跳跃。

\textbf{第 5 步（端点）。}$r=a$ 时每项 $J_0(x_n)=0$，级数为零；而 $f(a)=0$，两者一致，不存在跳跃。这与 \S8.3.2.5 的 $f=1$ 情形不同：那里 $f(a)=1$ 而级数给 $0$，端点必然跳跃，级数收敛到 $1/2$。

\answerstart{P20}{自编，$\mu$ 的判据}{四个柱坐标问题，判断 $\mu$ 的符号。}

判据只有一句：\textbf{哪一组边界是齐次的，就由哪一组定本征值；齐次条件要求对应的函数反复回到零，所以那一组必须是振荡型。}

\textbf{(a) 侧面 $u(a,\varphi,z)=0$（齐次），上下底给定（非齐次）。}侧面齐次 $\Rightarrow$ 径向函数要在 $r=a$ 为零 $\Rightarrow$ 贝塞尔，$\mu>0$；$Z$ 用 $e^{\pm\sqrt\mu z}$，其系数由上下底的非齐次数据定。

\textbf{(b) 上下底 $u=0$（齐次），侧面给定（非齐次）。}底面齐次 $\Rightarrow Z$ 要在 $z=0,h$ 为零 $\Rightarrow Z=\sin(n\pi z/h)$，$\mu=-\nu^2<0$，$\nu_n=n\pi/h$；径向方程是修正贝塞尔，$R=C I_m(\nu r)+D K_m(\nu r)$，实心柱丢 $K_m$。侧面条件用来定 $C$。

\textbf{(c) 无限长圆柱，无 $z$ 边界。}$z$ 方向没有边界条件，$Z$ 不能振荡也不能被某个 $z$ 条件定；取 $\mu=0$ 不损失一般性（$Z=A+Bz$，若问题与 $z$ 无关则 $Z$ 为常数）。这退化为 \S8.2 的二维极坐标问题：$\mu=0$，径向是欧拉方程。若题目明确 $u$ 与 $z$ 无关，就直接按二维做。

\textbf{(d) 侧面绝热（第二类齐次），上下底给定，轴对称。}$\varphi$ 无关 $\Rightarrow m=0$；侧面的第二类齐次条件 $\partial_ru(a)=0$ 对 $R=C_1+C_2\ln r$ 给 $C_2=0$，$R=C_1$ 自动满足；$Z=A+Bz$。这是 \S8.3.4.3 的 $\mu=0$ 分支。答案是一维线性温度 $u=A+Bz$，与半径无关——绝热侧面 + 轴对称的稳态，径向确实没有热流。

\aside{这道题的价值在于把“$\mu$ 的符号”从记公式变成判据：先问哪一组边界齐次，再问那一组要求的函数能不能反复回零，答案就唯一了。考试里若题目没有明说哪组齐次，就把所有的边界条件读一遍——非齐次的那一组是“数据入口”，齐次的那一组才是“谱的来源”。}

